% 推前测度：(a)点1,2,3,4权重0.1,0.2,0.3,0.4，T(1)=T(3)=a，T(2)=T(4)=b，得0.4delta_a+0.6delta_b；(b)lambda|[0,1]经T(x)=floor(2x)推前得1/2delta_0+1/2delta_1。
\begin{tikzpicture}[x=1cm,y=1cm,>=Stealth,
  v/.style={circle,draw=FigureStroke,fill=FigureBlueFill,line width=.7pt,minimum size=5.2mm,inner sep=0pt,font=\fontsize{8}{9}\selectfont},
  w/.style={circle,draw=FigureStroke,fill=FigureGreenFill,line width=.7pt,minimum size=6.2mm,inner sep=0pt,font=\fontsize{8}{9}\selectfont},
  ar/.style={draw=FigureMuted,line width=.6pt,->,shorten >=1pt,shorten <=1pt},
  note/.style={font=\FigureNoteFont\fontsize{7.5}{9}\selectfont,text=FigureMuted},
  tl/.style={font=\fontsize{7.5}{9}\selectfont}]
\begin{scope}
  \foreach \i/\y/\w in {1/1.2/0.1,2/0.4/0.2,3/-0.4/0.3,4/-1.2/0.4}{\node[v] (x\i) at (0,\y) {$\i$};\node[tl,anchor=east] at (-0.35,\y) {$\w$};}
  \node[w] (a) at (2.2,0.6) {$a$};
  \node[w] (b) at (2.2,-0.6) {$b$};
  \draw[ar] (x1)--(a); \draw[ar] (x3)--(a);
  \draw[ar] (x2)--(b); \draw[ar] (x4)--(b);
  \node[tl,anchor=west] at (2.55,0.6) {$0.4$};
  \node[tl,anchor=west] at (2.55,-0.6) {$0.6$};
  \node[note] at (1.1,1.75) {$T$：多对一};
  \node at (1.1,-2.0) {(a) 有限点的合并};
\end{scope}
\begin{scope}[xshift=4.5cm,xscale=0.8]
  \draw[FigureStroke,line width=.6pt] (-0.2,1.0)--(4.7,1.0);
  \draw[FigureBlue,line width=3pt] (0,1.0)--(2,1.0);
  \draw[FigureGreen,line width=3pt] (2,1.0)--(4,1.0);
  \foreach \x/\l in {0/0,2/\tfrac12,4/1}{\draw[FigureStroke,line width=.6pt] (\x,0.92)--(\x,1.08);\node[tl,above] at (\x,1.08) {$\l$};}
  \node[note,anchor=west] at (4.75,1.0) {$\lambda|_{[0,1]}$};
  \draw[FigureStroke,line width=.6pt] (-0.2,-0.9)--(4.7,-0.9);
  \foreach \x/\l in {0.5/0,3.5/1}{\draw[FigureStroke,line width=.6pt] (\x,-0.98)--(\x,-0.82);\node[tl,below] at (\x,-0.98) {$\l$};}
  \fill[FigureBlue] (0.4,-0.9) rectangle (0.6,-0.15);
  \fill[FigureGreen] (3.4,-0.9) rectangle (3.6,-0.15);
  \node[tl,above] at (0.5,-0.15) {$\tfrac12$};
  \node[tl,above] at (3.5,-0.15) {$\tfrac12$};
  \draw[ar] (1,0.8) to[bend right=10] (0.62,0.35);
  \draw[ar] (3,0.8) to[bend left=10] (3.38,0.35);
  \node[note,anchor=west] at (4.75,-0.9) {$T_\#\lambda|_{[0,1]}$};
  \node[note] at (2,0.35) {$T(x)=\lfloor2x\rfloor$};
  \node at (2,-2.0) {(b) 区间质量合并为点质量};
\end{scope}
\end{tikzpicture}
